\chapter{Rule Sheets}
\label{ch:rule-sheets}
\label{app:rules}

\section{Classical natural deduction}
\label{sec:rule-sheets-01}

Contexts in this sheet are finite sets of formulas.  A displayed rule is a rule
schema: uniformly substituting formulas for its formula metavariables produces
a rule instance.  In the propositional portion, \(\derives\) abbreviates
\(\derives_{\ND_{\CPL}}\); in the quantifier and identity portion it
abbreviates \(\derives_{\ND_{\FOL}}\).  The latter calculus contains every
rule of the former.  Assumptions enter derivations through
\[
\Rule[assumption]{}{\Gamma\cup\set{\varphi}\derives\varphi}.
\]

\subsection*{Conjunction}
\[
\Rule[\(\land I\)]{\Gamma\derives\varphi\qquad\Delta\derives\psi}
{\Gamma\cup\Delta\derives\varphi\land\psi}
\qquad
\Rule[\(\land E_1\)]{\Gamma\derives\varphi\land\psi}
{\Gamma\derives\varphi}
\qquad
\Rule[\(\land E_2\)]{\Gamma\derives\varphi\land\psi}
{\Gamma\derives\psi}.
\]

\subsection*{Disjunction}
\[
\Rule[\(\lor I_1\)]{\Gamma\derives\varphi}
{\Gamma\derives\varphi\lor\psi}
\qquad
\Rule[\(\lor I_2\)]{\Gamma\derives\psi}
{\Gamma\derives\varphi\lor\psi}
\]
\[
\Rule[\(\lor E\)]{
\Gamma\derives\varphi\lor\psi\qquad
\Delta\cup\set\varphi\derives\chi\qquad
\Sigma\cup\set\psi\derives\chi}
{\Gamma\cup\Delta\cup\Sigma\derives\chi}.
\]
The displayed occurrences of \(\varphi\) and \(\psi\) in the minor
derivations are discharged.

\subsection*{Implication}
\[
\Rule[\(\to I\)]{\Gamma\cup\set\varphi\derives\psi}
{\Gamma\derives\varphi\to\psi}
\qquad
\Rule[\(\to E\)]{\Gamma\derives\varphi\to\psi\qquad
\Delta\derives\varphi}
{\Gamma\cup\Delta\derives\psi}.
\]
The displayed occurrence of \(\varphi\) in implication introduction is
discharged.

\subsection*{Primitive negation, the bottom formula, and RAA}
\[
\Rule[\(\neg I\)]{\Gamma\cup\set{\varphi}\derives\bottom}
{\Gamma\derives\neg\varphi}
\qquad
\Rule[\(\neg E\)]{\Gamma\derives\varphi\qquad
\Delta\derives\neg\varphi}
{\Gamma\cup\Delta\derives\bottom}
\]
\[
\Rule[\(\bottom E\)]{\Gamma\derives\bottom}
{\Gamma\derives\varphi}
\qquad
\Rule[RAA]{\Gamma\cup\set{\neg\varphi}\derives\bottom}
{\Gamma\derives\varphi}.
\]
The rules \(\to I\), \(\neg I\), and RAA discharge the displayed assumption.
Negation is primitive in the active classical language.  Removing RAA gives
the primitive-negation presentation of its intuitionistic fragment.

\subsection*{Quantifiers}
\[
\Rule[\(\forall I\)]{\Gamma\derives\varphi}
{\Gamma\derives\forall x\,\varphi}
\side{\(x\notin\FV(\Gamma)\)}
\]
\[
\Rule[\(\forall E\)]{\Gamma\derives\forall x\,\varphi}
{\Gamma\derives\varphi[t/x]}
\]
\[
\Rule[\(\exists I\)]{\Gamma\derives\varphi[t/x]}
{\Gamma\derives\exists x\,\varphi}
\]
\[
\Rule[\(\exists E\)]{
\Gamma\derives\exists x\,\varphi\qquad
\Delta\cup\set{\varphi[y/x]}\derives\psi}
{\Gamma\cup\Delta\derives\psi}.
\]
\noindent\emph{Side condition.}  \(y\) is absent from \(\psi\), from every
other open assumption, and from \(\exists x\,\varphi\).
The displayed assumption \(\varphi[y/x]\) is discharged.  In the contextual
display, ``every other open assumption'' means every formula in
\(\Gamma\cup\Delta\); the side condition is therefore literally the one used
by the active natural-deduction calculus.

\subsection*{Alpha-equivalence}
\[
\Rule[\(\alpha\)]{\Gamma\derives\varphi}{\Gamma\derives\psi}
\side{\(\varphi\equiv_\alpha\psi\)}.
\]
This is the explicit bound-variable conversion rule of the active calculus;
it does not permit replacement by an arbitrary logically equivalent
formula.

\subsection*{Logical identity}
\[
\Rule[\(=I\)]{}{\derives t=t}
\qquad
\Rule[\(=E\)]{\Gamma\derives r=t\qquad\Delta\derives\varphi[r/x]}
{\Gamma\cup\Delta\derives\varphi[t/x]}.
\]
Identity elimination is symmetric because symmetry is derivable from these
rules.  Every substitution displayed in this sheet is the total
capture-avoiding substitution fixed in the syntax chapter, so no separate
substitution side condition is imposed.

\section{Classical sequent calculus}
\label{sec:rule-sheets-02}

A sequent is \(\Gamma\Rightarrow\Delta\), with finite multisets on both sides.
It is interpreted as
\(\bigwedge\Gamma\to\bigvee\Delta\), with the empty conjunction interpreted as
\(\topv\) and empty disjunction as \(\bottom\).

The propositional rules below are exactly the active calculus
\(\LK_{\CPL}\).  The subsequent quantifier and identity rules extend it to
\(\LK_{\FOL}\).

\subsection*{Initial sequents and cut}
\[
\Rule[id]{}{\Gamma,\varphi\Rightarrow\varphi,\Delta}
\qquad
\Rule[\(\bottom L\)]{}{\Gamma,\bottom\Rightarrow\Delta}
\]
\[
\Rule[cut]{\Gamma\Rightarrow\Delta,\varphi\qquad
\Sigma,\varphi\Rightarrow\Pi}
{\Gamma,\Sigma\Rightarrow\Delta,\Pi}.
\]

\subsection*{Structural rules}
\[
\Rule[WL]{\Gamma\Rightarrow\Delta}{\varphi,\Gamma\Rightarrow\Delta}
\quad
\Rule[WR]{\Gamma\Rightarrow\Delta}{\Gamma\Rightarrow\Delta,\varphi}
\]
\[
\Rule[CL]{\varphi,\varphi,\Gamma\Rightarrow\Delta}
{\varphi,\Gamma\Rightarrow\Delta}
\quad
\Rule[CR]{\Gamma\Rightarrow\Delta,\varphi,\varphi}
{\Gamma\Rightarrow\Delta,\varphi}.
\]
Exchange is built into the multiset convention.

\subsection*{Logical rules}
\[
\Rule[\(\neg L\)]{\Gamma\Rightarrow\Delta,\varphi}
{\Gamma,\neg\varphi\Rightarrow\Delta}
\qquad
\Rule[\(\neg R\)]{\Gamma,\varphi\Rightarrow\Delta}
{\Gamma\Rightarrow\Delta,\neg\varphi}
\]
\[
\Rule[\(\land L\)]{\varphi,\psi,\Gamma\Rightarrow\Delta}
{\varphi\land\psi,\Gamma\Rightarrow\Delta}
\qquad
\Rule[\(\land R\)]{\Gamma\Rightarrow\Delta,\varphi\qquad
\Gamma\Rightarrow\Delta,\psi}
{\Gamma\Rightarrow\Delta,\varphi\land\psi}
\]
\[
\Rule[\(\lor L\)]{\varphi,\Gamma\Rightarrow\Delta\qquad
\psi,\Gamma\Rightarrow\Delta}
{\varphi\lor\psi,\Gamma\Rightarrow\Delta}
\qquad
\Rule[\(\lor R\)]{\Gamma\Rightarrow\Delta,\varphi,\psi}
{\Gamma\Rightarrow\Delta,\varphi\lor\psi}
\]
\[
\Rule[\(\to L\)]{\Gamma\Rightarrow\Delta,\varphi\qquad
\Gamma,\psi\Rightarrow\Delta}
{\Gamma,\varphi\to\psi\Rightarrow\Delta}
\qquad
\Rule[\(\to R\)]{\varphi,\Gamma\Rightarrow\Delta,\psi}
{\Gamma\Rightarrow\Delta,\varphi\to\psi}.
\]

\subsection*{Quantifiers}
\[
\Rule[\(\forall L\)]{\Gamma,\varphi[t/x]\Rightarrow\Delta}
{\Gamma,\forall x\,\varphi\Rightarrow\Delta}
\qquad
\Rule[\(\exists R\)]{\Gamma\Rightarrow\Delta,\varphi[t/x]}
{\Gamma\Rightarrow\Delta,\exists x\,\varphi}
\]
\[
\Rule[\(\forall R\)]{\Gamma\Rightarrow\Delta,\varphi[y/x]}
{\Gamma\Rightarrow\Delta,\forall x\,\varphi}
\qquad
\Rule[\(\exists L\)]{\Gamma,\varphi[y/x]\Rightarrow\Delta}
{\Gamma,\exists x\,\varphi\Rightarrow\Delta}.
\]
In \(\forall R\), the eigenvariable \(y\) is free in no formula of
\(\Gamma,\Delta,\forall x\,\varphi\).  In \(\exists L\), it is free in no
formula of \(\Gamma,\Delta,\exists x\,\varphi\).  Every displayed
substitution is total and capture-avoiding.

\subsection*{Logical identity}
\[
\Rule[\(=I\)]{}{\Rightarrow t=t}
\qquad
\Rule[\(=L\)]{\Gamma,\varphi[r/x]\Rightarrow\Delta}
{\Gamma,r=t,\varphi[t/x]\Rightarrow\Delta}.
\]
The symmetric replacement instance is derivable.  These displays, together
with the propositional rules, are the complete active first-order sequent
calculus used in this reference.

\section[Modal calculus for K]{Modal calculus for \(\K\)}
\label{sec:rule-sheets-03}

The Hilbert calculus \(\K\) contains all substitution instances of classical
classically valid propositional formulas, the axiom schema
\[
  \Box(\varphi\to\psi)\to(\Box\varphi\to\Box\psi),
\]
modus ponens, and necessitation:
\[
  \Rule[nec]{\derives_{\K}\varphi}{\derives_{\K}\Box\varphi}.
\]
Necessitation applies to theorems, not to an arbitrary formula depending on an
undischarged assumption.  Adding \(\Box\varphi\to\varphi\),
\(\Box\varphi\to\Box\Box\varphi\), and
\(\Diamond\varphi\to\Box\Diamond\varphi\) yields the standard systems
\(\mathsf T\), \(\Sfour\), and \(\Sfive\) in the combinations defined in the
modal chapter.

